% GENERATED FILE. Edit canonical lessons, then run npm run generate:books.
% canonical-source-hash: fnv1a64:db30ed802e6a2eaa
% canonical-lessons: PT-C464-geada, PT-C464-orvalho, PT-C464-congelar, PT-C464-congelador, PT-C464-derreter, PT-R464-revisao-457-464

\chapter{O gelo}
\label{ch:pt-o-gelo}
\hlchaptermodality{464}

\begin{chapteropening}
\textbf{By the end of this chapter:} \emph{I can describe a cold spell: the frost, the dew, freezing, the freezer and melting.}

Complete the last lesson of chapter 464: I can describe a cold spell: the frost, the dew, freezing, the freezer and melting.
\end{chapteropening}

\section[a geada]{a geada — the frost}

\label{lesson:PT-C464-geada}



\begin{quote}
Before the new one: say the Portuguese for hail, then the Portuguese for the downpour.
\end{quote}

\subsection*{You'll want to know: a geada}
\textbf{a geada} — \textquotedblleft{}the frost\textquotedblright{}.

\begin{quote}
De manhã, a relva estava branca de geada.
\end{quote}

\begin{quote}
\emph{In the morning, the grass was white with frost.}
\end{quote}

\emph{Cair geada} — to frost over.

\begin{cousinweb}
\emph{geada} is inherited from Latin \emph{gelāta} 'frozen' (EtymDB), from \emph{gelāre} 'to freeze', on \emph{gelū} 'frost'.
\end{cousinweb}

\begin{grammarlens}[title={What you've built}]
One more word for telling, explaining and describing at B1.
\end{grammarlens}

\subsection*{Your turn}
\begin{itemize}
\raggedright
  \item \emph{Say it:} \emph{a geada}
  \item \emph{Say it:} \emph{a geada}, once more
  \item \emph{Say it:} a sentence of your own with \emph{a geada}
  \item \emph{Recall:} say the Portuguese for hail, then the Portuguese for the downpour, then say \emph{a geada} again
\end{itemize}

\begin{tcolorbox}[breakable,colback=teal!4,colframe=teal!35!black,title={Before you move on}]
What does a geada mean? (\textquotedblleft{}The frost\textquotedblright{}.) Say it once more.
\end{tcolorbox}
\section[o orvalho]{o orvalho — the dew}

\label{lesson:PT-C464-orvalho}



\begin{quote}
Before the new one: say the Portuguese for the downpour, then the Portuguese for the frost.
\end{quote}

\subsection*{You'll want to know: o orvalho}
\textbf{o orvalho} — \textquotedblleft{}the dew\textquotedblright{}.

\begin{quote}
O orvalho molha os pés de manhã cedo.
\end{quote}

\begin{quote}
\emph{The dew wets your feet early in the morning.}
\end{quote}

Drops on the grass at dawn.

\begin{cousinweb}
\emph{orvalho} is usually traced to Latin \emph{rōs, rōris} 'dew' (Lewis \& Short), but the details are not settled.
\end{cousinweb}

\begin{grammarlens}[title={What you've built}]
One more word for telling, explaining and describing at B1.
\end{grammarlens}

\subsection*{Your turn}
\begin{itemize}
\raggedright
  \item \emph{Say it:} \emph{o orvalho}
  \item \emph{Say it:} \emph{o orvalho}, once more
  \item \emph{Say it:} a sentence of your own with \emph{o orvalho}
  \item \emph{Recall:} say the Portuguese for the downpour, then the Portuguese for the frost, then say \emph{o orvalho} again
\end{itemize}

\begin{tcolorbox}[breakable,colback=teal!4,colframe=teal!35!black,title={Before you move on}]
What does o orvalho mean? (\textquotedblleft{}The dew\textquotedblright{}.) Say it once more.
\end{tcolorbox}
\section[congelar]{congelar — to freeze}

\label{lesson:PT-C464-congelar}



\begin{quote}
Before the new one: say the Portuguese for the frost, then the Portuguese for the dew.
\end{quote}

\subsection*{You'll want to know: congelar}
\textbf{congelar} — \textquotedblleft{}to freeze\textquotedblright{}.

\begin{quote}
No inverno, a água do lago congela.
\end{quote}

\begin{quote}
\emph{In winter, the water in the lake freezes.}
\end{quote}

Also food: \emph{congelar a sopa}.

\begin{cousinweb}
\emph{congelar} is a learned borrowing from Latin \emph{congelāre} 'to freeze together' (EtymDB), on \emph{gelū} 'frost'.
\end{cousinweb}

\begin{grammarlens}[title={What you've built}]
One more word for telling, explaining and describing at B1.
\end{grammarlens}

\subsection*{Your turn}
\begin{itemize}
\raggedright
  \item \emph{Say it:} \emph{congelar}
  \item \emph{Say it:} \emph{congelar}, once more
  \item \emph{Say it:} a sentence of your own with \emph{congelar}
  \item \emph{Recall:} say the Portuguese for the frost, then the Portuguese for the dew, then say \emph{congelar} again
\end{itemize}

\begin{tcolorbox}[breakable,colback=teal!4,colframe=teal!35!black,title={Before you move on}]
What does congelar mean? (\textquotedblleft{}To freeze\textquotedblright{}.) Say it once more.
\end{tcolorbox}
\section[o congelador]{o congelador — the freezer}

\label{lesson:PT-C464-congelador}



\begin{quote}
Before the new one: say the Portuguese for the dew, then the Portuguese for to freeze.
\end{quote}

\subsection*{You'll want to know: o congelador}
\textbf{o congelador} — \textquotedblleft{}the freezer\textquotedblright{}.

\begin{quote}
Há peixe no congelador.
\end{quote}

\begin{quote}
\emph{There is fish in the freezer.}
\end{quote}

Part of \emph{o frigorífico} or a separate box.

\begin{cousinweb}
\emph{congelador} is built on \emph{congelar} 'to freeze', a learned borrowing from Latin \emph{congelāre} (EtymDB).
\end{cousinweb}

\begin{grammarlens}[title={What you've built}]
One more word for telling, explaining and describing at B1.
\end{grammarlens}

\subsection*{Your turn}
\begin{itemize}
\raggedright
  \item \emph{Say it:} \emph{o congelador}
  \item \emph{Say it:} \emph{o congelador}, once more
  \item \emph{Say it:} a sentence of your own with \emph{o congelador}
  \item \emph{Recall:} say the Portuguese for the dew, then the Portuguese for to freeze, then say \emph{o congelador} again
\end{itemize}

\begin{tcolorbox}[breakable,colback=teal!4,colframe=teal!35!black,title={Before you move on}]
What does o congelador mean? (\textquotedblleft{}The freezer\textquotedblright{}.) Say it once more.
\end{tcolorbox}
\section[derreter]{derreter — to melt}

\label{lesson:PT-C464-derreter}



\begin{quote}
Before the new one: say the Portuguese for to freeze, then the Portuguese for the freezer.
\end{quote}

\subsection*{You'll want to know: derreter}
\textbf{derreter} — \textquotedblleft{}to melt\textquotedblright{}.

\begin{quote}
Com o sol, a neve derrete depressa.
\end{quote}

\begin{quote}
\emph{In the sun, the snow melts quickly.}
\end{quote}

Ice cream \emph{derrete} too.

\begin{cousinweb}
\emph{derreter} is usually traced to Latin \emph{dēterere} 'to rub away' (Lewis \& Short), but the origin is not settled; EtymDB has no entry for it.
\end{cousinweb}

\begin{grammarlens}[title={What you've built}]
One more word for telling, explaining and describing at B1.
\end{grammarlens}

\subsection*{Your turn}
\begin{itemize}
\raggedright
  \item \emph{Say it:} \emph{derreter}
  \item \emph{Say it:} \emph{derreter}, once more
  \item \emph{Say it:} a sentence of your own with \emph{derreter}
  \item \emph{Recall:} say the Portuguese for to freeze, then the Portuguese for the freezer, then say \emph{derreter} again
\end{itemize}

\begin{tcolorbox}[breakable,colback=teal!4,colframe=teal!35!black,title={Before you move on}]
What does derreter mean? (\textquotedblleft{}To melt\textquotedblright{}.) Say it once more.
\end{tcolorbox}
\section[(dialogue)]{Revisão dos capítulos 457–464}

\label{lesson:PT-R464-revisao-457-464}



\begin{quote}
Say the Portuguese for the oak and the pine.
\end{quote}

\subsection*{Your turn}
Say these aloud:

\begin{itemize}
\raggedright
  \item the oak, the pine, the cork oak, cork, to prune
  \item the mountain range, the hillside, the summit, the gorge, to climb
  \item the big rock, the tide, the seashell, to tan, to dive
  \item the wild boar, the deer, the hare, the falcon, to graze
  \item the lizard, the viper, the toad, the snail, to crawl
  \item muggy, the drought, the light breeze, sunstroke, to sweat
  \item the thunderstorm, to thunder, to flash with lightning, hail, the downpour
  \item the frost, the dew, to freeze, the freezer, to melt
\end{itemize}

\begin{tcolorbox}[breakable,colback=teal!4,colframe=teal!35!black,title={Before you move on}]
The oak? (\textbf{o carvalho}.) The pine? (\textbf{o pinheiro}.) The cork oak? (\textbf{o sobreiro}.) Cork? (\textbf{a cortiça}.) To prune? (\textbf{podar}.) The mountain range? (\textbf{a serra}.) The hillside? (\textbf{a encosta}.) The summit? (\textbf{o cume}.) The gorge? (\textbf{o desfiladeiro}.) To climb? (\textbf{escalar}.) The big rock? (\textbf{o rochedo}.) The tide? (\textbf{a maré}.) The seashell? (\textbf{a concha}.) To tan? (\textbf{bronzear}.) To dive? (\textbf{mergulhar}.) The wild boar? (\textbf{o javali}.) The deer? (\textbf{o veado}.) The hare? (\textbf{a lebre}.) The falcon? (\textbf{o falcão}.) To graze? (\textbf{pastar}.) The lizard? (\textbf{o lagarto}.) The viper? (\textbf{a víbora}.) The toad? (\textbf{o sapo}.) The snail? (\textbf{o caracol}.) To crawl? (\textbf{rastejar}.) Muggy? (\textbf{abafado}.) The drought? (\textbf{a seca}.) The light breeze? (\textbf{a aragem}.) Sunstroke? (\textbf{a insolação}.) To sweat? (\textbf{transpirar}.) The thunderstorm? (\textbf{a trovoada}.) To thunder? (\textbf{trovejar}.) To flash with lightning? (\textbf{relampejar}.) Hail? (\textbf{o granizo}.) The downpour? (\textbf{o aguaceiro}.) The frost? (\textbf{a geada}.) The dew? (\textbf{o orvalho}.) To freeze? (\textbf{congelar}.) The freezer? (\textbf{o congelador}.) To melt? (\textbf{derreter}.)
\end{tcolorbox}
